% reputation.tex -- Reputation System and Social Memory

\section{Reputation and Social Memory}
\label{sec:reputation}

The reputation system introduces \emph{cross-episode social memory} to Kant,
enabling agents to build, communicate, and act on opponent reputations across
multiple interactions. This tests whether LLMs can reason about social
reputation and adapt to opponent track records---a capability central to
sustained cooperation in repeated social
settings~\citep{nowak2005evolution,axelrod1984evolution}.

\subsection{Gossip Variant}
\label{sec:gossip}

The gossip variant follows the composable \texttt{apply\_*} pattern established
by cheap talk and rule proposal (Sections~\ref{sec:mechanism}
and~\ref{sec:governance}). For base actions $\{A, B\}$ and ratings
$\{$\texttt{trustworthy}, \texttt{untrustworthy}, \texttt{neutral}$\}$,
gossip produces the expanded action space:
\[
\bigl\{\texttt{gossip\_}\langle\text{rating}\rangle\texttt{\_}\langle\text{action}\rangle
\;\big|\; \text{rating} \in R,\; \text{action} \in A_{\text{base}}\bigr\}
\]
yielding $|R| \times |A_{\text{base}}|$ composite actions (e.g., six for PD).
Crucially, payoffs depend \emph{only} on the base action component---gossip is
non-binding cheap talk about reputation, not a payoff-relevant signal:
\[
u_i(\texttt{gossip\_}r\texttt{\_}a,\; \texttt{gossip\_}r'\texttt{\_}b) = u_i(a, b)
\]

Three pre-registered gossip games are provided: Gossip Prisoner's Dilemma,
Gossip Stag Hunt, and Gossip Hawk-Dove. The variant composes freely with exit,
cheap talk, and rule proposal, enabling multi-layered communication studies.

\subsection{Memory Store}
\label{sec:memory-store}

The \texttt{CogneeMemoryStore} provides persistent cross-episode memory backed
by the Cognee knowledge graph engine. After each episode, a structured text
summary is ingested:

\begin{quote}\small\ttfamily
Game Interaction Report\\
Agent: agent\_0 | Opponent: tit\_for\_tat | Game: prisoners\_dilemma\\
Rounds: 10 | Agent Score: 30 | Opponent Score: 30\\
Cooperation Rate: 1.0\\
Actions: R1: cooperate vs cooperate; R2: cooperate vs cooperate; \ldots
\end{quote}

\noindent The knowledge graph is rebuilt after each ingestion
(\texttt{cognify}), enabling semantic queries about opponent behavior patterns.
When Cognee is unavailable, the system degrades gracefully to fast in-memory
statistics tracking cooperation rates and interaction counts.

\subsection{Reputation Environment}
\label{sec:rep-env}

The \texttt{ReputationEnvironment} wraps \texttt{KantEnvironment} to
inject reputation data into observations and record interactions:

\begin{enumerate}[nosep]
    \item \textbf{Before each episode}: query the memory store for opponent
          reputation and inject it into \texttt{obs.metadata}.
    \item \textbf{During play}: extract gossip ratings from composite actions
          and record them in the knowledge graph.
    \item \textbf{After each episode}: record the full episode summary with
          cooperation rate and scores.
\end{enumerate}

\noindent This creates a feedback loop where agent behavior in early episodes
shapes opponent reputation, which in turn influences agent reasoning in
subsequent episodes. The architecture tests whether LLMs can integrate
historical reputation signals into strategic decision-making---a core
requirement for trustworthy multi-agent deployment.
